\documentclass[10pt,a4paper]{amsart}
\usepackage[margin=19mm]{geometry}
\usepackage{amsmath,amssymb,mathtools}
\usepackage[scr=rsfs,cal=euler]{mathalfa}
\usepackage{xfrac}
\usepackage[unicode,hidelinks,bookmarks=false]{hyperref}
\newcommand{\ie}{i.e.\ }
\newcommand{\cf}{cf.\ }
\newcommand{\resp}{respectively\ }
\newcommand{\Subsection}{\subsection}
\providecommand{\nobreakdash}{\nobreak\hbox{-}}
\setlength{\emergencystretch}{2em}
\multlinegap0pt
\usepackage{amssymb}
\usepackage{mathtools}
\usepackage{xcolor}
\newtheorem{thm}{Theorem}
\newtheorem{lem}{Lemma}
\newtheorem{propo}{Proposition}
\newcommand{\Bcal}{\mathcal{B}}
\newcommand{\C}{\mathbb{C}}
\newcommand{\Cfp}[1]{C^{\ast}_{\mathrm{fp}}({#1})}
\newcommand{\CfpCartan}[1]{\ell^{\infty}_{\mathrm{fp}}({#1})}
\newcommand{\Ecal}{\mathcal{E}}
\newcommand{\Fcal}{\mathcal{F}}
\newcommand{\Hcal}{\mathcal{H}}
\newcommand{\N}{\mathbb{N}}
\newcommand{\abs}[1]{{\left\lvert #1 \right \rvert}}
\newcommand{\asdim}[1]{\mathrm{asdim}({#1})}
\newcommand{\bop}[1]{\Bcal({#1})}
\newcommand{\ddimn}[3]{\dim_{\mathrm{diag}}({#1}\subseteq{#2}\, ;~{#3})}
\DeclareMathOperator{\dist}{d}
\newcommand{\distX}[2]{{\dist_X ( #1, #2 )}}
\newcommand{\e}{\varepsilon}
\newcommand{\gUnMatImg}[3]{{ \mathsf{g}_{#3}^{(#1), #2}}}
\newcommand{\id}[1]{\mathrm{Id}_{#1}}
\newcommand{\indF}{\mathbf{1}}
\newcommand{\myhyp}{\textnormal{-}\hskip0pt\relax}
\newcommand{\norm}[1]{{\left\lVert #1 \right \rVert}}
\newcommand{\normBig}[1]{{\Big\lVert #1 \Big \rVert}}
\newcommand{\normbig}[1]{{\big\lVert #1 \big \rVert}}
\newcommand{\prn}[1]{{(#1)}}
\newcommand{\rAr}{\longrightarrow}
\newcommand{\unMatImg}[3]{{\Ecal_{#3}^{(#1), #2}}}
\title{Generalised Diagonal Dimension and applications to large-scale geometry}
\author{Christos Kitsios}
\date{Source: arXiv:2604.07237v1 (CC BY 4.0)}
\begin{document}
\maketitle

\noindent\emph{Source and licence.} Generalised Diagonal Dimension and applications to large-scale geometry. Version arXiv:2604.07237v1 (CC BY 4.0), distributed by arXiv under the Creative Commons Attribution 4.0 International licence, http://creativecommons.org/licenses/by/4.0/, as recorded in the arXiv licence metadata for this exact version and shown on the abstract page https://arxiv.org/abs/2604.07237v1. Full licence notice: \texttt{licenses/arxiv-2604.07237-CC-BY-4.0.txt}.\par\smallskip

\noindent\textit{Editorial scope.} Counted unit: the article's thm thm:lowerBoundFPAlg (source0.tex lines 1421--1427). The article's complete original proof of it is delivered with it (source0.tex lines 2044--2259, one \texttt{proof} environment of 216 lines, which the article places away from the statement). No mathematics is written, repaired or re-derived: the statement and the proof are the article's own lines, in the article's own notation, and no retained line is substituted or re-flowed.  Retained uncounted as the source-local context the proof needs: lines 1912--1919; lines 2007--2029.  Nothing was omitted from any retained span, so the package carries no omission record.  No retained line contains a citation, so the wrapper carries no bibliography and no bibliography entry.  No retained line contains a figure environment, an image-inclusion command or drawing code, so no image is delivered and the package declares no asset dependency.  Editorial formatting only, in addition: an AMS \texttt{amsart} preamble that restates the article's own declarations and macros used by the retained lines, namely the environments \texttt{thm}, \texttt{lem}, \texttt{propo} on separate counters, together with the standard packages those lines need.  The retained environments print in this document as Theorem 1, Lemma 1, Propo 1 in order of appearance, whereas the article prints the same items with the numbering of its own full document.  The complete native source of the article as distributed in the arXiv e-print for this exact version is bundled unchanged as \texttt{sources/198078/source0.tex}.
    \begin{lem}\label{lem:LowerBoundProof1}
    If
    \begin{equation*}
        S\coloneq \max_{i=0,\hdots,d} \bigg(\sup_{x \in U^{(i)}} \big{\lvert}{ [ x ]_{i}}   \big{\rvert} \bigg) < \infty ,
    \end{equation*}
    then
    \[\asdim{X} \le d.\]
    \end{lem}

    \begin{propo}  \label{prop:capCpcMaps}
    The  maps 
    \[ A \overset{\widehat{\psi}}{\longrightarrow} \widehat{F} \overset{\widehat{\phi}}{\longrightarrow} A,\]
    are
    completely positive contractive maps
    such that:
        \begin{enumerate} \renewcommand{\labelenumi}{(\roman{enumi})}
            \item for each $a \in A$, it holds
            \begin{equation*}
             \phi \circ \psi (a) = \frac{1}{1+\e^2/81} \widehat{\phi} \circ \widehat{\psi} (a),
            \end{equation*}
            \item 
                 for each $a \in \Fcal \cup \Fcal^2$, it holds
            \begin{equation*}
        \normBig{\widehat{\phi} \circ \widehat{\psi}(a) - a} < \e^2/27,
    \end{equation*}
            \item for each $a \in \Fcal$, $b \in \widehat{F}$, with $\norm{b} \le 1$, it holds
            \begin{equation*} \label{ineq:capCpcMapsBound}
        \normBig{ \widehat{\phi} \big( \widehat{\psi}(a) b \big) - \widehat{\phi}\big(\widehat{\psi}(a)\big) \widehat{\phi}\big(b\big)} 
        < 6 \Big(\frac{\e^2}{81}\Big)^{1/2}  < \e.
    \end{equation*}
        \end{enumerate}
    \end{propo}

\begin{thm}  \label{thm:lowerBoundFPAlg}
    Let $X$ be a uniformly locally finite metric space.
    Then
    \begin{align*}
         \ddimn{\CfpCartan{X}}{\Cfp{\Hcal_X}}{\ell^\infty(\N, \bop{\Hcal})} \ge \asdim{X}.
    \end{align*}
\end{thm}

    \begin{proof}[Proof of Theorem~\ref{thm:lowerBoundFPAlg}]
    By Lemma~\ref{lem:LowerBoundProof1}, in order to prove the theorem, it suffices to show that 
     \begin{equation*}
        S = \max_{i=0,\hdots,d} \bigg(\sup_{x \in U^{(i)}} \big{\lvert}{ [ w ]_{i}} \big{\rvert} \bigg) < \infty .
    \end{equation*}

        Fix ${i \in \{0,\hdots, d\}}$.
      Let ${w \in U^{\prn{i}}}$ and without loss of generality assume that ${w \in U^{(i),j_0}_{k_0}}$,
      for some ${j_0}$ and ${k_0}$. 
        Suppose that ${\widetilde{w} \in U^{\prn{i}}}$ is such that ${w \sim_i \widetilde{w}}$. 
        Then there exist 
        \[{w_1, \hdots, w_\kappa \in U^{\prn{i}}}\]
        with 
    \begin{itemize}
        \item $w_1 = w,$
        \item $w_\kappa = \widetilde{w},$
        \item ${\distX{w_l}{w_{l+1}} \le r}$ for each ${l \in \{1,\hdots,\kappa-1\}}$.
    \end{itemize}

    Fix ${l \in \{1,\hdots,\kappa-1\}}$ and set 
    ${z=w_l}$ and ${z'=w_{l+1}}$. 
    Without loss of generality, assume that
    ${z \in U^{(i),j}_k}$ and ${z' \in U^{(i),j'}_{k'}}$ for some  ${j, j'}$ and ${k, k'}$. 
        Set 
        \begin{equation*}
            x \coloneq {{\overline{\sigma}}_{k,1}^{(i),j}}(z)  \in U^{(i),j}_1 
        \end{equation*}
        and
        \begin{equation*}
        y \coloneq {\overline{\sigma}_{k',1}^{(i),j'}}(z')  \in U^{(i),j'}_1.
        \end{equation*}
       Note that
        \[{{\overline{\sigma}}_{1,k}^{(i),j}} \big(x \big) = {{\overline{\sigma}}_{1,k}^{(i),j}} \big({{\overline{\sigma}}_{k,1}^{(i),j}}(z) \big) = z\]
        and
        \[{\overline{\sigma}_{1,k'}^{(i),j'}}\big(y \big) = {\overline{\sigma}_{1,k'}^{(i),j'}}\big({\overline{\sigma}_{k',1}^{(i),j'}}(z') \big)
        = z'. \]
        
        We will show that 
        $j = j'$ and $x=y$, and, in particular, it holds
        \begin{equation*}
            z' = {\overline{\sigma}_{1,k'}^{(i),j}}\big({{\overline{\sigma}}_{k,1}^{(i),j}}(z) \big).
        \end{equation*}

        We define the operators
        \begin{equation*}
            T_x \coloneq \indF_x \, \gUnMatImg{i}{j}{1,k} \, \indF_{z} \, \gUnMatImg{i}{j}{k,1} \, \indF_x \in D_A
        \end{equation*}
        and 
        \begin{equation*}
        \widetilde{T}_{z} \coloneq \indF_{z} \, \gUnMatImg{i}{j}{k,1} T_x \gUnMatImg{i}{j}{1,k} \, \indF_{z} \in D_A,
        \end{equation*}
        which are supported on $\{x\}$ and $\{z\}$, respectively.
        Since ${z \in U^{(i),j}_k}$, it holds
         \begin{equation*}
        \normbig{\unMatImg{i}{j}{k,k}  \widetilde{T}_{z} \unMatImg{i}{j}{k,k} }  > \eta^2.
    \end{equation*}
         Moreover, using the above estimate and the definition of ${\widehat{\phi}}$ we see that
        \begin{equation*}
            \norm{ \widehat{\phi}\Big(e_{k,k}^{(i),j} \Big)  \widetilde{T}_{z} \unMatImg{i}{j}{k,k}}  >  \delta^2/2.
        \end{equation*}
        Set 
       \begin{equation*}
           \widehat{T} \coloneq  \widehat{\phi}\left(e_{k,k}^{(i),j} \right)   \widetilde{T}_{z} \unMatImg{i}{j}{k,k} \in D_A,
       \end{equation*}
       and observe that $\widehat{T}$ is supported on $\{z\}$, 
       since $\phi$ maps diagonal elements into $D_A$.
        
        Since ${\distX{z'}{z} \le r}$, there exists a unique ${m \in \{1,\hdots,M\}}$ such that
     \[{a_m(z',z) = \id{\Hcal}}.\] 
        Then, using the matrix representation of operators in $A$, 
    it is straightforward to see that
    \begin{equation*}
                \widehat{T}(z,z) 
                = (a_m \widehat{T} a_m^\ast) (z',z') = (\indF_{z'} a_m \widehat{T} a_m^\ast) (z',z'),
    \end{equation*}
    where ${(\indF_{z'} a_m \widehat{T} a_m^\ast) (z',z')}$ is the ${(z',z')}$\myhyp{}entry in the matrix representation of the operator ${\indF_{z'} a_m \widehat{T} a_m^\ast\in A}$.

        By combining the above
        \begin{align*}
          \frac{\delta^2}{2}
          & <  \normBig{\widehat{T}} 
          = \normBig{\widehat{T}(z,z)}
          = \normBig{(\indF_{z'} a_m \widehat{T} a_m^\ast) (z',z')}
          \le \normBig{\indF_{z'} a_m \widehat{T} a_m^\ast }.
        \end{align*}

        Since ${a_m \in \Fcal}$, using Proposition~\ref{prop:capCpcMaps}.(ii), we obtain
         \begin{align*}
          \frac{\delta^2}{2}
          < \normBig{\indF_{z'} \widehat{\phi} \big(\widehat{\psi}(a_m) \big)  \widehat{T} a_m^\ast } + \frac{\e^2}{27}.
        \end{align*}
        Then, by the definition of $\widehat{T}$ and Proposition~\ref{prop:capCpcMaps}.(iii), we have
        \begin{align*}
          \frac{\delta^2}{2}
           & < \normBig{\indF_{z'} \widehat{\phi} \big(\widehat{\psi}(a_m) \big)  \widehat{T} a_m^\ast } + \frac{\e^2}{27}\\
           & = \normBig{\indF_{z'}  \, \widehat{\phi} \big(\widehat{\psi}(a_m) \big) \widehat{\phi}\big(e_{k,k}^{(i),j} \big)  \widetilde{T}_{z} \unMatImg{i}{j}{k,k} a_m^\ast} + \frac{\e^2}{27}\\
           & \le \normBig{\indF_{z'} \, \widehat{\phi}\big(\widehat{\psi}(a_m) e_{k,k}^{(i),j} \big)   \widetilde{T}_{z} \unMatImg{i}{j}{k,k} a_m^\ast} + \frac{\e^2}{27} +\e,
        \end{align*}
        and, using the definition of
        $\widehat{\psi}$ and $\widehat{\phi}$,
        we obtain
        \begin{align*}
          \frac{\delta^2}{2}
           & \le   \normBig{\indF_{z'} \, {\phi}\big({\psi}(a_m) e_{k,k}^{(i),j} \big)   \widetilde{T}_{z} \unMatImg{i}{j}{k,k} a_m^\ast} + \frac{\e^2}{27} +\e.
        \end{align*}
  
        Define 
        \[{h^{(i)} \coloneq \phi^{(i)} \big( 1_{F^{(i)} \otimes B} \big)}\]
        and suppose that
         \[\pi^{(i)} \colon F^{(i)} \otimes B~\rAr~ \bop{\ell^2(X,\Hcal)}\]
         is a supporting ${\ast}$\myhyp{}homomorphism of the 
         order zero map $\phi^{(i)}$, that is,
        \[{\phi^{(i)}(b) = h^{(i)} \pi^{(i)}(b)},\]
        for each ${b \in F^{(i)} \otimes B}$.
    Condition~(6) in the definition of generalised diagonal dimension implies that
    \begin{equation*}
        \indF_{z' }= \pi^{(i)}\big(e^{{(i)},{j'}}_{k',k'} \big) \; \indF_{z' } \;\pi^{(i)}\big(e^{{(i)},{j'}}_{k',k'} \big) .
    \end{equation*}

        Note that
        ${\unMatImg{i}{j'}{1,k'} \indF_{z'} \unMatImg{i}{j'}{k',1}}$
        is supported in $\{y\}$,
        since ${z' = \overline{\sigma}_{1,k'}^{(i),j'}\big(y \big)}$.
        Then, using the definition of ${\unMatImg{i}{j'}{1,k'}}$
    and the definition of $\pi^{(i)}$, we obtain that the operator
    \begin{equation*}
        P_y \coloneq \pi^{(i)}\big(e^{\prn{i},j'}_{1,k'} \big) \; \indF_{z'} \;\pi^{(i)}\big(e^{\prn{i},j'}_{k',1} \big) \in D_A 
    \end{equation*}
    is supported in $\{y\}$.
    We have
    \begin{align*}
        \pi^{(i)}\big(e^{\prn{i},j'}_{k',1} \big)  P_y \pi^{(i)}\big(e^{\prn{i},j'}_{1,k'} \big)
        = \pi^{(i)}\big(e^{\prn{i},j'}_{k',k'} \big) \; \indF_{z' } \;\pi^{(i)}\big(e^{\prn{i},j'}_{k',k'} \big) 
        = 
        \indF_{z'}.
    \end{align*}

        By replacing the above in the previous inequality, we obtain
        \begin{align} \label{ineq:xNEyIneq}
          \frac{\delta^2}{2}
           <  \normBig{\pi^{(i)}\big(e^{\prn{i},j'}_{k',1} \big)  P_y \pi^{(i)}\big(e^{\prn{i},j'}_{1,k'} \big) {\phi}\big({\psi}(a_m) e_{k,k}^{(i),j} \big)   \widetilde{T}_{z} \unMatImg{i}{j}{k,k} a_m^\ast} + \frac{\e^2}{27} +\e.
        \end{align}

        Since $\phi^{(i)}$ is an order zero map with $\pi^{\prn{i}}$ a supporting $\ast$\myhyp{}homomorphism,
        ${j \ne j'}$ 
        implies that the first summand above is zero, and then
        \[\frac{\delta^2}{2}
          < \frac{\e^2}{27} +\e,\]
        which is a contradiction. 
        Therefore, $j=j'$.

        Assume now that ${x \ne y}$ and
        set 
        \[v \coloneq e^{{(i)},{j}}_{1,k'} \psi(a_m) e^{{(i)},{j}}_{k,1} \in F^{\prn{i}} \otimes B.\] 
        Observe that, by its definition,
        $v$ is of the form ${v = e^{{(i)},{j},\C}_{1,1} \otimes b}$ for some ${b \in B}$.
        
        Then, 
        using that ${\pi^{\prn{i}}}$ is a supporting $\ast$\myhyp{}homomorphism of $\phi^{(i)}$, 
        and that 
        \[\gUnMatImg{i}{j}{k,1} =g_\delta\big(h^{(i)} \big) \pi^{(i)}\big(e^{{(i)},{j}}_{k,1}\big),\]
        we obtain
        \begin{align*}
            \pi^{(i)}\big(e^{\prn{i},j}_{1,k'} \big) {\phi}\big({\psi}(a_m) e_{k,k}^{(i),j} \big)   \widetilde{T}_{z}
            & = \pi^{(i)}\big(e^{\prn{i},j}_{1,k'} \big) {\phi}\big({\psi}(a_m) e_{k,k}^{(i),j} \big)    \gUnMatImg{i}{j}{k,1} T_x \gUnMatImg{i}{j}{1,k}  \indF_{z}\\
            & =  {\phi}\big( e^{\prn{i},j}_{1,k'} {\psi}(a_m) e_{k,k}^{(i),j} e^{\prn{i},j}_{k,1} \big)   g_\delta\big(h^{(i)}\big) T_x \gUnMatImg{i}{j}{1,k}  \indF_{z}\\
           &  =  {\phi}\big( v \big)   g_\delta\big(h^{(i)}\big) T_x \gUnMatImg{i}{j}{1,k}  \indF_{z}.
        \end{align*}
    
        The following:
        \begin{itemize}
            \item ${v \in D^{\prn{i}} \otimes B}$,
            \item ${h^{(i)} = \phi^{(i)} \big( 1_{F^{(i)} \otimes B} \big)}$,
            \item $\phi$ maps the diagonal into $D_A$,
        \end{itemize}
        imply that 
         \[{{\phi}\big( v \big)   g_\delta\big(h^{(i)}\big) \in D_A}.\]
        Recall that ${P_y \in D_A}$ is supported on $\{y\}$ and ${T_x \in D_A}$ is supported on $\{x\}$. 
        Then, using that ${x \ne y}$, we obtain
        \begin{align*}
            P_y \pi^{(i)}\big(e^{\prn{i},j}_{1,k'} \big) {\phi}\big({\psi}(a_m) e_{k,k}^{(i),j} \big)   \widetilde{T}_{z}
            = P_y {\phi}\big( v \big)   g_\delta\big(h^{(i)}\big) T_x \gUnMatImg{i}{j}{1,k}  \indF_{z} = 0.
        \end{align*}
          
         Therefore, by combining the above and
         Inequality~(\ref{ineq:xNEyIneq}), we have
         \begin{align*}
          \frac{\delta^2}{2}
           & <  \normBig{\pi^{(i)}\big(e^{\prn{i},j'}_{k',1} \big)  P_y \pi^{(i)}\big(e^{\prn{i},j'}_{1,k'} \big) {\phi}\big({\psi}(a_m) e_{k,k}^{(i),j} \big)   \widetilde{T}_{z} \unMatImg{i}{j}{k,k} a_m^\ast} + \frac{\e^2}{27} +\e
           = \frac{\e^2}{27} +\e,
        \end{align*}
        which is a contradiction. 
        Thus, $x=y$.
  
         Therefore, we obtain
         \begin{equation} \label{eq:LowerBoundProof1}
                w_{l+1} = z' 
         ={\overline{\sigma}_{1,k'}^{(i),j}}\big(y \big)
        = {\overline{\sigma}_{1,k'}^{(i),j}}\big(x \big)
        = {\overline{\sigma}_{1,k'}^{(i),j}}\big({{\overline{\sigma}}_{k,1}^{(i),j}}(z) \big)
        = {\overline{\sigma}_{1,k'}^{(i),j}}\big({{\overline{\sigma}}_{k,1}^{(i),j}}(w_l) \big). 
         \end{equation}

    Recall that ${w = w_1, w_2, \hdots, \widetilde{w}=w_\kappa \in U^{\prn{i}}}$ are such that ${\distX{w_l}{w_{l+1}} \le r}$ for each ${l}$. 
    Then, by applying the previous (Equation~\eqref{eq:LowerBoundProof1}) recursively to the pairs ${(w_{m}, w_{m+1})}$ for each ${m=1,\hdots,\kappa-1}$, 
    we obtain 
    \begin{equation*}
        \widetilde{w} = w_\kappa \in \Big\{{{\overline{\sigma}}_{1,k'}^{(i),j_0}}\big({\overline{\sigma}}_{k_0,1}^{(i),j_0}(w) \big) \colon k' = 1,\hdots,s^{(i),j_0}  \Big\}.
    \end{equation*}
    This implies 
        \[ \big{\lvert}{ [ w ]_{i}} \big{\rvert}  \le s^{(i),j_0} \]
        and, therefore,
        \begin{equation*}
            S = \max_{i=0,\hdots,d} \bigg(\sup_{x \in U^{(i)}} \big{\lvert}{ [ w ]_{i}} \big{\rvert} \bigg) \le \max_{ i=0,\hdots,d  } \bigg(\sup_{j=1,\hdots,r^{(i)}} \abs{ s^{(i),j} }\bigg)  < \infty . \qedhere
        \end{equation*}
\end{proof}

\end{document}
